\chapter{Última dispersión y bariogénesis}

\section{Mirar hacia el exterior es mirar hacia el pasado}

A distintas distancias vemos el universo en distintos momentos. Para los objetos cercanos, el tiempo retrospectivo no es muy importante. El universo de hace un año, o de hace cien años, no era muy distinto del actual. Pero, si miramos lo bastante lejos, vemos un universo más joven. Un universo más joven es más denso, y uno más denso es más caliente.

A cierta distancia suficientemente grande, el universo estaba tan caliente como la superficie del Sol. No lo vemos así porque los fotones se han desplazado al rojo en su trayecto hasta aquí. Un factor de desplazamiento al rojo de aproximadamente mil implica una temperatura observada menor por ese mismo factor:
\[
T_{\mathrm{obs}} \sim \frac{T_{\mathrm{emit}}}{1+z},
\qquad
1+z \sim 10^3.
\]\footnote{Esta expresión es una transcripción editorial concisa del argumento verbal sobre el desplazamiento al rojo expuesto en 00:02:20--00:02:39.}

Ahí es donde el universo se volvió opaco. Estaba compuesto principalmente de hidrógeno, con algo de helio. Por encima de unos pocos miles de grados, el hidrógeno está ionizado. Así que, más allá de esa superficie, no hay átomos neutros, sino iones: principalmente protones y electrones que se mueven libremente. Allí los fotones también tienen más energía; son, básicamente, rayos X. Si se forma un átomo por casualidad, las colisiones o un fotón lo vuelven a ionizar enseguida.

Un gas ionizado es un conductor eléctrico, y un conductor no deja pasar la radiación sin alterarla. Las cargas responden a los campos eléctricos, oscilan de un lado a otro y vuelven a emitir radiación. Por eso, más allá de esa distancia, el universo es opaco para nosotros. La superficie de última dispersión es como la superficie visible del Sol: no es una pared material, sino un lugar donde el medio pasa de manera bastante repentina de opaco a transparente.

Allí, los fotones superan en número a las demás partículas por un factor de aproximadamente \(10^8\). Lo mismo ocurre hoy, aunque actualmente los fotones tienen tan poca energía que no constituyen una parte importante de la densidad de energía presente. Cuando la temperatura cae por debajo de la temperatura de ionización del hidrógeno, el universo se vuelve transparente y podemos ver a través de él.

La observación directa termina ahí. La siguiente pregunta es hasta dónde más atrás pueden llevarnos la física atómica, nuclear y de partículas conocidas.

\section{Remontarse aún más en el tiempo}

Justo después de la última dispersión, además de fotones, las partículas principales son protones y electrones, junto con algunos núcleos de helio y unos pocos deuterones. Prácticamente no hay antiprotones ni positrones. Pero, si nos remontamos más atrás, la temperatura acaba siendo lo bastante alta para que dos fotones choquen y produzcan un electrón y un positrón:
\[
\gamma + \gamma \leftrightarrow e^- + e^+.
\]

El umbral básico es que los fotones deben aportar al menos la energía en reposo del par:
\[
E \ge 2m_e c^2.
\]
Por debajo de ese umbral, la luz ordinaria no basta. Por encima, pueden crearse pares electrón--positrón, y también puede producirse la reacción inversa. Una olla de materia a esa temperatura contendría reacciones entre partículas que cumplirían el papel que desempeñan las reacciones químicas en la materia ordinaria.

Justo por encima del umbral sigue habiendo más fotones que electrones y positrones. La masa en reposo del electrón es costosa de producir. Pero, cuando la temperatura sube muy por encima del umbral, la energía en reposo del electrón pasa a ser menos importante que la energía cinética, y las densidades se vuelven comparables. Entonces, electrones, positrones y fotones tienen aproximadamente la misma densidad y chocan entre sí a toda velocidad.

\begin{figure}[ht]
\centering
\includegraphics[width=0.62\textwidth]{lecture_06_figure_02.png}
\caption{Esquema del umbral de producción de pares que compara las escalas de masa del electrón y del protón.}
\lectureframe{00:18:47}
\end{figure}

La siguiente pregunta es qué ocurre con los protones. Ya están presentes y, para entonces, son muy densos, pero todavía no hay suficiente energía térmica para producir pares protón--antiprotón en abundancia. La razón es sencilla:
\[
m_p \approx 2000\,m_e,
\qquad
E \ge 2m_p c^2
\]
es un umbral mucho más alto.\editorialnote{En la imagen, la etiqueta de menor masa es menos nítida que la superior \(2m_e c^2\), pero la comparación verbal se establece explícitamente entre los umbrales del electrón y del protón.}

Esto ya es mucho más caliente que el interior de una estrella, donde las antipartículas no están presentes en ninguna abundancia considerable.

Pero incluso esa descripción es solo provisional. Cuando la temperatura alcanza el valor necesario para producir protones y antiprotones libremente, también es lo bastante alta para descomponer los protones en quarks. Así que la sopa realmente caliente se describe mejor como una mezcla de electrones, positrones, quarks, antiquarks, fotones, neutrinos, gluones y otras especies. Una reacción esquemática sencilla de la pizarra es
\[
e^- + e^+ \to \gamma \to q + \bar q.
\]

\begin{classroomqa}
\lecturetimestamp{00:20:05}
\audiencequestion{¿Se trata de mesones o de prequarks?}
\lecturerresponse{Los mesones están hechos de quarks. Cuando la temperatura es lo bastante alta para desintegrar un protón, también desintegra los mesones.}
\end{classroomqa}

Hasta ahora, el razonamiento ha seguido el tiempo hacia atrás. Ahora invirtamos el sentido y preguntemos qué ocurre cuando la sopa caliente y simétrica se enfría.

\section{Neutralidad eléctrica y un universo caliente simétrico}

A temperatura suficientemente alta, el universo contiene enormes densidades de partículas y antipartículas. Los electrones y positrones aparecen en pares, los quarks y antiquarks aparecen en pares, y la imagen natural es que las abundancias de cada partícula y su antipartícula son casi iguales.

La neutralidad eléctrica requiere cierta cautela. El universo es eléctricamente neutro hoy. Un recurso mnemotécnico útil es imaginar, no nuestro universo real, sino uno cerrado y acotado, como una esfera tridimensional. Pensemos primero en una esfera bidimensional como modelo. Las líneas de campo eléctrico solo pueden empezar o terminar en una carga. Una carga positiva emite líneas de fuerza; una negativa las absorbe. En una esfera cerrada, esas líneas no pueden escapar sin más hacia el infinito. Vuelven a converger en el lado opuesto. Así que, como recurso mnemotécnico, un universo cerrado y acotado no puede tener carga neta. La cuestión no es si nuestro universo real es literalmente cerrado. La cuestión es que las cargas positivas y negativas a gran escala parecen cancelarse con extraordinaria precisión.

La neutralidad eléctrica \emph{no} significa que el número de protones sea igual al de antiprotones. Significa que la carga positiva total es igual a la carga negativa total. Los protones, los positrones y todas las demás especies con carga positiva, en conjunto, equilibran los electrones y todas las demás especies con carga negativa, también en conjunto.

Hasta donde podemos confiar en la física conocida, y quizá un poco más allá, el universo era lo bastante caliente para mantener partículas y antipartículas en abundancia. Cualquier experimento de acelerador que se nos ocurra tuvo un análogo en alguna etapa del universo temprano. La conjetura natural es que el estado inicial fuera casi simétrico: igual número de electrones y positrones, igual número de quarks y antiquarks y, si se toman los protones y antiprotones como representantes de los bariones y antibariones,
\[
N_p = N_{\bar p}.
\]

Esa conjetura resulta atractiva porque las leyes de la física son casi simétricas entre partículas y antipartículas. Un positrón se comporta de manera muy parecida a un electrón; no exactamente igual, pero casi. Esa simetría aproximada es el problema.

\section{Hacer avanzar la película}

Supongamos que el universo comienza en un estado caliente con abundancias iguales o casi iguales de partículas y antipartículas y luego se enfría.

Mientras la temperatura sea lo bastante alta, las colisiones energéticas pueden producir protones y antiprotones, que también pueden aniquilarse y convertirse de nuevo en fotones. Pero, tarde o temprano, la temperatura cae por debajo del umbral del protón. Entonces, los pares protón--antiprotón aún pueden aniquilarse y convertirse en fotones, pero, tras una pequeña expansión, esos fotones han perdido suficiente energía para ya no poder volver a producir pares protón--antiprotón. Es un camino de sentido único. Lo mismo ocurre más tarde con los electrones y positrones, a una temperatura mucho menor que el umbral del protón, pero aún muy superior a las temperaturas estelares.

Como la expansión es lenta en comparación con las tasas de reacción microscópicas, durante este proceso el universo se mantiene, en su mayor parte, cerca del equilibrio térmico. Esto es lo que aquí se entiende por expansión adiabática. En ese contexto, la afirmación más precisa es que la entropía cambia muy poco. Para una sopa caliente de partículas, la entropía equivale esencialmente al número de partículas, salvo por factores de orden unidad. Así que los protones, antiprotones, electrones y positrones que se aniquilan son reemplazados, aproximadamente, por fotones.

Si el número de protones y antiprotones fuera exactamente igual, todos se aniquilarían. Lo mismo les ocurriría a los electrones y positrones. Solo quedarían fotones. Pero si hubiera aunque fuera un pequeño exceso de protones sobre antiprotones, los pares correspondientes se aniquilarían y el exceso restante perduraría. Un ejemplo concreto lo aclara: si hubiera cien protones y cincuenta antiprotones, los cincuenta antiprotones encontrarían cincuenta protones, se aniquilarían con ellos y quedarían cincuenta protones. La neutralidad de carga exige entonces que también quede un exceso de electrones sobre positrones.

Esto permite deducir la asimetría temprana a partir del recuento actual de fotones. Hoy hay aproximadamente entre \(10^9\) y \(10^{10}\) veces más fotones que protones o electrones:
\[
\frac{N_\gamma}{N_p} \sim 10^9\text{--}10^{10}.
\]
Así que en la sopa caliente temprana había apenas un protón excedente por cada mil millones de pares protón--antiprotón. Ese minúsculo desequilibrio es la razón por la que todavía existen átomos.

\begin{classroomqa}
\lecturetimestamp{00:42:25}
\audiencequestion{¿Por qué no interpretar el exceso estadísticamente, como algo que se genera durante el enfriamiento con una probabilidad muy pequeña?}
\lecturerresponse{Para obtener un protón excedente a partir de un estado inicialmente simétrico, habría que crear un protón sin crear también un antiprotón. Además, habría que crear con él una carga negativa, por ejemplo, un electrón. El modelo estándar conocido no contiene ningún proceso de baja energía de ese tipo. A las energías estudiadas, se conserva el número bariónico, así que el enfriamiento ordinario no genera un exceso neto.}
\end{classroomqa}

Así que la carga eléctrica no es el obstáculo. Un ejemplo esquemático sería
\[
\gamma + \gamma \to p + e^-.
\]\footnote{Esta expresión es una transcripción editorial del ejemplo verbal presentado en 00:44:16--00:46:35. Su propósito es aislar la idea de que la carga eléctrica podría conservarse aunque se violara el número bariónico.}
En conjunto, esto sería eléctricamente neutro. Pero viola otra regla empírica. En el modelo estándar sometido a prueba no se puede crear un protón sin crear también un antiprotón. Para lo que nos ocupa, la cantidad conservada relevante es el número bariónico,
\[
B=(N_p+N_n)-(N_{\bar p}+N_{\bar n}).
\]\footnote{Esta fórmula concisa formaliza la definición verbal dada en 00:45:18--00:45:40.}
A las energías conocidas, esa cantidad se conserva, así que la asimetría no puede generarse durante la etapa tardía y bien comprobada de la historia de enfriamiento.

\section{Bariogénesis}

Este es el tema conocido como bariogénesis. Supongamos que el universo realmente comenzó con igual número de bariones y antibariones. ¿Qué tendrían que permitir las leyes de la física para inclinar ese estado simétrico hacia el pequeño exceso observado?

En la pizarra, esto se redujo a dos expresiones simbólicas,
\[
N_p=N_{\bar p},
\qquad
N_p-N_{\bar p},
\]
y luego a una lista de condiciones que conectan la primera con la segunda.

\begin{figure}[ht]
\centering
\includegraphics[width=0.72\textwidth]{lecture_06_figure_04.png}
\caption{Condiciones para generar una asimetría bariónica a partir de un estado inicialmente simétrico.}
\lectureframe{01:22:26}
\end{figure}

\editorialnote{En esta imagen, las palabras de la lista numerada solo se distinguen parcialmente. Los tres puntos que siguen corresponden a la explicación oral que los acompaña.}

La primera condición es la violación del número bariónico. Si el número bariónico es una ley de conservación exacta, un universo que comienza con número bariónico neto nulo permanece así.

La segunda condición es un sesgo entre partículas y antipartículas. La violación del número bariónico por sí sola no basta. Si la naturaleza permite un proceso que produce protones y su proceso reflejado, que produce antiprotones, con la misma probabilidad, el desequilibrio medio sigue siendo cero. Es necesario que las tasas sean distintas.

Las fluctuaciones estadísticas ordinarias también son demasiado pequeñas. Si el número total de partículas es aproximadamente \(10^{90}\), incluidos los fotones, una fluctuación del orden de \(\sqrt{N}\) da apenas \(10^{45}\), muy lejos de los aproximadamente \(10^{80}\) protones del universo.

La tercera condición es que, en algún momento, el universo haya estado suficientemente alejado del equilibrio. Aquí el argumento recurre al teorema CPT. Si un gas caliente ha permanecido siempre en equilibrio, una película reproducida hacia atrás se ve igual que una reproducida hacia adelante. En ese contexto, la simetría entre partículas y antipartículas se restaura efectivamente mediante la simetría exacta que combina conjugación de carga, paridad e inversión temporal. Para que quede fijado un sesgo real, la historia efectiva debe verse distinta al avanzar y al retroceder en escalas de tiempo microscópicas. En una etapa suficientemente temprana, la expansión del universo era lo bastante rápida para que esto ocurriera.

Así pues, las tres condiciones son:
\begin{enumerate}
\item violación del número bariónico;
\item violación de la simetría entre partículas y antipartículas, más precisamente, violación de CP;
\item una etapa fuera del equilibrio térmico.
\end{enumerate}

\begin{classroomqa}
\lecturetimestamp{01:09:08}
\audiencequestion{Al relacionarlas con CPT, ¿la condición 2 es realmente CP y la condición 3 se refiere realmente a la cuestión de la dirección temporal?}
\lecturerresponse{Sí. La condición 2 debería formularse como violación de CP. La condición 3 es el hecho de que la historia efectiva estuviera fuera del equilibrio; no es una ley de conservación independiente ni equivale a violar el número bariónico.}
\end{classroomqa}

Con estas tres condiciones, incluso un universo que comienza con igual número de protones y antiprotones puede inclinarse en un sentido o en el otro. Lo que sigue siendo incierto es si la naturaleza realmente viola el número bariónico y, de ser así, cómo.

\section{Lo que se sabe y lo que no}

La parte aún abierta de la historia es la violación del número bariónico. Nunca hemos observado directamente en la naturaleza una reacción que cambie
\[
(N_p+N_n)-(N_{\bar p}+N_{\bar n}).
\]
Eso quizá solo signifique que las reacciones pertinentes requieren temperaturas mucho más altas que cualquiera de las que se han estudiado directamente en el laboratorio, o que son extraordinariamente improbables a bajas energías. Pero, por ahora, esta sigue siendo la incógnita.

Por eso la desintegración del protón es tan importante en este debate. Si el número bariónico no es exacto, en principio el protón debería ser inestable. Sin embargo, es extraordinariamente estable. El límite inferior citado aquí es del orden de
\[
\tau_p \gtrsim 10^{34}\ \text{years}.
\]
Este resultado procede de enormes experimentos que utilizan volúmenes muy grandes de agua equipados con detectores de luz y enterrados a gran profundidad. No demuestra que no haya violación del número bariónico. Solo significa que cualquier violación de ese tipo debe estar extraordinariamente suprimida a bajas energías.

En el debate también aparece una observación aparte sobre la física conocida. Alrededor de \(100\) a \(250\,\mathrm{GeV}\), las interacciones débiles se parecen mucho más a las interacciones electromagnéticas. Eso podría importar en modelos detallados, pero no es el paso decisivo del argumento que se presenta aquí.

La materia oscura no cambia mucho la lógica básica. En la mayoría de las teorías interactúa tan débilmente con la materia ordinaria que, para esta cuestión, es esencialmente un espectador.

\begin{classroomqa}
\lecturetimestamp{01:24:49}
\audiencequestion{¿Cómo afecta la gravedad a todo esto?}
\lecturerresponse{Principalmente de dos maneras. Primero, sin expansión cósmica no habría ninguna etapa fuera del equilibrio. Segundo, la evaporación de los agujeros negros ofrece un motivo de peso para sospechar que el número bariónico no puede ser una ley de conservación exacta para siempre. Fuera de eso, las interacciones gravitatorias ordinarias entre partículas son demasiado débiles para importar mucho aquí.}
\end{classroomqa}

Entonces surge una pregunta relacionada con la carga eléctrica. ¿El mismo argumento de los agujeros negros también pondría en entredicho la conservación de la carga? La respuesta es no. Un agujero negro cargado tiene líneas de flujo eléctrico que se extienden hacia el exterior, y esas líneas dan cuenta de la carga. Ese agujero negro no se evapora por completo, sino que se aproxima a un remanente extremal. Así que el argumento de los agujeros negros pone en entredicho el número bariónico, no la carga eléctrica.

Otra pregunta plantea si la sopa de quarks muy densa del universo temprano debería considerarse simplemente un agujero negro. La respuesta es no. La cuestión no es una masa densa y estática. El universo estaba en expansión. El argumento de los agujeros negros es un experimento mental sobre el destino del número bariónico, no una descripción literal del universo temprano como una región atrapada.

\begin{classroomqa}
\lecturetimestamp{01:31:14}
\audiencequestion{¿Se sabría cómo modificar el modelo estándar para incluir ese tipo de reacciones?}
\lecturerresponse{El problema es el contrario. Es fácil añadir demasiada violación del número bariónico y hacer que el protón se desintegre con demasiada rapidez. El verdadero problema es mantener el protón tan estable como es y, a la vez, permitir los procesos que se necesitan en el universo temprano.}
\end{classroomqa}

\section{La entropía y las últimas preguntas sobre los neutrones}

Volvamos ahora a hacer el balance. En un universo que se expande adiabáticamente, la entropía se conserva aproximadamente. Para una sopa caliente de partículas elementales, la entropía equivale esencialmente al número de partículas, salvo por factores de orden unidad. Por eso, la discusión no requiere una enorme producción tardía de fotones de la nada. Las aniquilaciones, en su mayor parte, reorganizan el contenido de la sopa y lo convierten en radiación.

La versión anterior del problema planteaba la historia al revés. En lugar de preguntar por qué sobrevivieron tan pocos protones en relación con los fotones, se preguntaba lo contrario: ¿por qué hay tantos fotones en comparación con los protones? Pero la pregunta más precisa no es por qué el universo generó tanta entropía en una etapa tardía. Es por qué la asimetría entre partículas y antipartículas era tan pequeña desde el principio: apenas un protón excedente por cada mil millones de pares protón--antiprotón.

Un neutrón libre sí se desintegra, pero ese proceso \emph{no} viola el número bariónico:
\[
n \to p + e^- + \bar\nu.
\]
El neutrón tiene el mismo número bariónico que el protón.

Para un neutrón libre, la energía disponible para la desintegración es muy pequeña, porque la masa del neutrón solo supera ligeramente la masa del estado final. Esa es una de las razones por las que la desintegración es tan lenta según los estándares de la física de partículas elementales. Cuando un neutrón está ligado en un núcleo, la energía de enlace puede eliminar por completo esa pequeña cantidad de energía disponible, y entonces el neutrón ligado puede ser estable.

\begin{classroomqa}
\lecturetimestamp{01:43:03}
\audiencequestion{¿Los neutrones simplemente desaparecen del universo o también se crean?}
\lecturerresponse{La primera respuesta es que los neutrones libres se pierden por desintegración. Pero esa no es toda la historia. Hay fuentes reales de producción de neutrones, incluidas las reacciones nucleares en las estrellas y la formación de estrellas de neutrones. Así que el número total de neutrones no disminuye simplemente en un solo sentido.}
\end{classroomqa}